\documentclass[11pt,openany]{report}
\usepackage[a4paper,margin=1in]{geometry}
\usepackage[T1]{fontenc}
\usepackage{lmodern}
\usepackage{hyperref}
\usepackage{fancyhdr}
\usepackage{titlesec}
\usepackage{longtable}
\usepackage{xcolor}

\hypersetup{
  colorlinks=true,
  linkcolor=blue,
  urlcolor=blue,
  pdftitle={kira2diffeq Manual},
  pdfauthor={kira2diffeq}
}

\definecolor{docblue}{RGB}{0,61,153}
\definecolor{docgray}{RGB}{245,245,245}

\titleformat{\chapter}{\normalfont\Large\bfseries\color{docblue}}{\thechapter.}{0.5em}{}
\titleformat{\section}{\normalfont\large\bfseries\color{docblue}}{\thesection}{0.5em}{}
\titleformat{\subsection}{\normalfont\bfseries}{\thesubsection}{0.5em}{}

\pagestyle{fancy}
\fancyhf{}
\fancyhead[L]{kira2diffeq Manual}
\fancyhead[R]{\thepage}
\renewcommand{\headrulewidth}{0.4pt}

\newcommand{\code}[1]{\texttt{#1}}

\begin{document}

\begin{titlepage}
  \begin{center}
    \vspace*{3cm}
    {\Huge\bfseries kira2diffeq Manual}\\[1em]
    {\large Differential equations from Kira reduction output}\\[2cm]
    \rule{12cm}{0.5pt}\\[1cm]
    {\large Project version: local repository guide}\\[0.8cm]
    {\normalsize\today}\\[2cm]
    \vfill
  \end{center}
\end{titlepage}

\tableofcontents
\clearpage

\chapter{Project scope}
kira2diffeq{} is a C++ utility that reads Kira-style reduction output and
constructs differential equations for master integrals.

It performs these operations:
\begin{enumerate}
  \item Read topology and kinematic configuration.
  \item Parse \code{kira2math} reductions and \code{masters.final}.
  \item Differentiate master integrals with respect to selected invariants.
  \item Reduce shifted integrals back to masters.
  \item Assemble and export DE matrices.
\end{enumerate}

This package does not perform diagram generation or IBP solving.

\chapter{Build and verification}

\section{Prerequisites}
\begin{itemize}
  \item CMake $\ge$ 3.16
  \item GiNaC
  \item yaml-cpp
  \item A C++17 compiler
\end{itemize}

\section{Standard commands}
\begin{verbatim}
cd /path/to/kira2diffeq
cmake -S . -B build
cmake --build build -j4
ctest --test-dir build --output-on-failure
\end{verbatim}

\section{Regression coverage}
\begin{itemize}
  \item \texttt{test\_sunrise}: parses config, reduction, differentiation, matrix assembly.
  \item \texttt{cli\_smoke}: end-to-end CLI run using local fixtures.
\end{itemize}

\chapter{Command line usage}

Executable: \code{build/kira2diffeq-cli}

\section{Options}
\begin{verbatim}
Usage: kira2diffeq-cli [options]
  -c <dir>   Config directory (integralfamilies.yaml, kinematics.yaml)
  -m <file>  kira2math .m output file
  -f <file>  masters.final file
  -t <name>  Topology name
  -v <var>   Kinematic variable (default: all)
  -o <fmt>   Output: text | math | all (default: text)
  -e         Use epsilon (d = 4 - 2*eps) in output
  -O <file>  Write output to file instead of stdout
  -h         Show this help
\end{verbatim}

\section{Example}
\begin{verbatim}
./build/kira2diffeq-cli \
  -c tests/fixtures/cli_smoke \
  -m tests/fixtures/cli_smoke/reduction.m \
  -f tests/fixtures/cli_smoke/masters.final \
  -t cli_smoke \
  -v t \
  -o text \
  -O /tmp/cli_smoke_de.txt
\end{verbatim}

Output file is created when \code{-O} is given.

\section{Output formats}
\begin{itemize}
  \item \code{text}: default, matrix form with $\mathrm{MI}_i$ notation.
  \item \code{math}: Mathematica syntax blocks.
  \item \code{all}: all variables in Mathematica format.
  \item \code{-e}: substitutes \code{d -> 4 - 2 eps} in output.
\end{itemize}

\section{Debug trace}
\code{differentiate\_master} is silent by default.
Enable detailed logs with:
\begin{verbatim}
export KIRA2DIFFEQ_DEBUG=1
./build/kira2diffeq-cli ...
\end{verbatim}

\chapter{Required inputs}

\section{Configuration directory}
\textbf{Required files:}
\begin{itemize}
  \item \code{kinematics.yaml}
  \item \code{integralfamilies.yaml}
\end{itemize}

\section{kinematics.yaml}
Expected structure:
\begin{itemize}
  \item incoming/outgoing momenta
  \item kinematic invariants
  \item scalarproduct rules
  \item optional momentum conservation
  \item optional \code{symbol\_to\_replace\_by\_one}
\end{itemize}

\section{integralfamilies.yaml}
Expected entries:
\begin{itemize}
  \item topology \code{name}
  \item \code{loop\_momenta}
  \item \code{top\_level\_sectors}
  \item \code{propagators}
\end{itemize}

Example snippets:
\begin{verbatim}
# standard propagator
- ["l1+p1-k2", t]

# bilinear ISP
- {bilinear: [["k1", "l1"], 0]}
\end{verbatim}

\section{masters.final and reduction file}
\code{masters.final} has one line per master, for example:
\begin{verbatim}
cli_smoke[1]
\end{verbatim}

\code{kira2math} lines map targets to linear sums of masters:
\begin{verbatim}
topo[1,2,0] -> +topo[0,1,0]*(1), +topo[1,1,0]*(-2)
\end{verbatim}

\chapter{Internal architecture}
\begin{itemize}
  \item \textbf{parser}: config and reduction data.
  \item \textbf{differentiator}: derivative rules + Gram-chain substitutions.
  \item \textbf{reducer}: map shifted integrals into master basis.
  \item \textbf{exporter}: text and Mathematica output.
\end{itemize}

Core files:
\begin{itemize}
  \item \code{src/config\_reader.cpp}
  \item \code{src/math\_parser.cpp}
  \item \code{src/differentiator.cpp}
  \item \code{src/reducer.cpp}
  \item \code{src/exporter.cpp}
  \item \code{src/main.cpp}
\end{itemize}

\chapter{Troubleshooting}

\section{Missing reduction entries}
Warning pattern:
\begin{verbatim}
Warning: integral X[...] not found in reduction table
\end{verbatim}

Interpretation: a shifted integral from differentiation is not reducible and not a master.

Fix list:
\begin{itemize}
  \item add the missing reduction rule to \code{.m}
  \item check topology name/ordering and propagator definitions
  \item verify invariant symbols between config and reduction are identical
\end{itemize}

\section{Parser errors}
\begin{itemize}
  \item check YAML structure in both config files
  \item ensure \code{kira2math} lines use the expected \code{target -> +master*(coeff)} form
  \item verify the selected topology with \code{-t}
\end{itemize}

\chapter{Project-context notes}
For a staged collinear workflow, keep regulated substitutions (for example \(p_{12}\))
in intermediate algebra and only apply strict limits after cancellations and checks.

\chapter{Quick project map}
\begin{itemize}
  \item \code{MANUAL.md}
  \item \code{MANUAL.tex}
  \item \code{CMakeLists.txt}
  \item \code{tests/CMakeLists.txt}
  \item \code{tests/fixtures/cli\_smoke/}
  \item \code{src/} and \code{include/}
\end{itemize}

\chapter{Appendix}
\section{Smoke workflow}
\begin{verbatim}
cmake --build build -j4
./build/kira2diffeq-cli \
  -c tests/fixtures/cli_smoke \
  -m tests/fixtures/cli_smoke/reduction.m \
  -f tests/fixtures/cli_smoke/masters.final \
  -t cli_smoke -v t -o text -O /tmp/cli_smoke_de.txt
\end{verbatim}

\end{document}
